\section{Comparison with Related Work}\label{sec:disc:comparison}

Classical \gls{elint} pipelines surveyed in \cref{sec:related:classical} deinterleave with \gls{toa} histograms or with density clustering on raw \gls{pdw} coordinates, then match trains against a library.
The encoder here replaces that first stage with a learned embedding, and it never performs the second.
Scores in \cref{ch:experiments} are partition agreement, not equipment names.
A raw-feature \gls{dbscan} control is the Feature \gls{dbscan} row of \cref{tab:exp:attn_em,tab:exp:attn_md}.
It oversplits emitters.
Vanilla is already well above that control, so the encoder is doing work that density clustering on the same scaled \glspl{pdw} does not.
On windows where frequency and pulse width collide, as in \cref{fig:exp:feat_w100_em,fig:exp:feat_maxmode}, a metric on those coordinates alone is not enough. The learned embeddings still separate the platforms.

Supervised radar classifiers in \cref{sec:related:dl_radar} assume a global emitter catalogue and a softmax at inference.
\textcite{jonsson2023supcon} trains a related contrastive encoder on the same simulator lineage, still as closed-set identification.
The present scores cannot be stacked against that accuracy. The tasks differ.
Emitter labels here are recording-local, and both branches are clustered rather than classified.
That is closer to \textcite{gunn2025deinterleaving}, who train a transformer for deinterleaving and cluster the embeddings with HDBSCAN.
The Vanilla row of \cref{tab:exp:attn_em} is the analogous existence test, on a different corpus and with \gls{dbscan} rather than HDBSCAN.
Gunn et al.\ stop at one partition. The second column of \cref{tab:exp:attn_md} is the addition, and \cref{tab:exp:ablation_loss} shows that it is not free: joint Vanilla deinterleaves less cleanly than an emitter-only copy of the same encoder.

The attention variants go one step past that counterpart.
Gunn et al.\ use standard dot-product attention and leave \gls{toa} as an input feature.
RoPE-TOA and Bias put time into the attention product itself, and on this test set they close the emitter tail that Vanilla leaves.
RoPE-az does not, once incidence is already in the token.
None of these variants is a published radar baseline. They are drop-in changes of the same encoder, which is the comparison RQ3 asks for.

Deep clustering methods in \cref{sec:related:deep_clustering} usually learn one partition, often with the number of clusters chosen in advance.
The two branches here are not a hierarchy of one relation. Modes are not subtypes of emitters.
A single embedding cannot keep both similarities, and \cref{tab:exp:ablation_loss} is the empirical trace of that conflict in the shared trunk.
The comparison with multi-task softmax heads in \cref{sec:related:joint} is the same point in other words. Those heads need globally consistent class names. The emitter identifiers in this corpus are not that.
